% \iffalse meta-comment
%
% measuredtex.dtx
%%
%% Permission is hereby granted, free of charge, to any person obtaining
%% a copy of this software and associated documentation files, to deal in
%% the Software without restriction, including without limitation the rights
%% to use, copy, modify, merge, publish, distribute, sublicense, and/or sell
%% copies of the Software, subject to the conditions stated in the MIT
%% License.
%%
%
% \fi
%
% \iffalse
%<*driver>
\documentclass{l3doc}
\EnableCrossrefs
\CodelineIndex
\RecordChanges
\begin{document}
  \DocInput{measuredtex.dtx}
\end{document}
%</driver>
% \fi
%
% \changes{v0.1}{2026/06/02}{Initial public version.}
%
% \GetFileInfo{measuredtex.sty}
%
% \title{The \textsf{measuredtex} package\thanks{%
%   This document corresponds to \textsf{measuredtex}~v0.1,
%   dated 2026-06-02.}}
% \author{Blair Hall\\Measurement Standards Laboratory of New Zealand}
% \date{2026-06-02}
%
% \maketitle
%
% \begin{abstract}
% The \textsf{measuredtex} package provides semantic markup for
% measurement-model documents. It supplies environments for typesetting
% model equations while writing structured model metadata to a sidecar
% file, commands for recording definitions of mathematical terms, and a
% lightweight mechanism for collecting textual concepts used in the
% document. The package currently writes model and mathematical-term records to
% \texttt{\cs{jobname}.models}. Textual concept records collected with
% \cs{mdconcept} are written at the end of the document to
% \texttt{\cs{jobname}.terms}.
% \end{abstract}
%
% \tableofcontents
%
% \section{Introduction}
%
% Measurement-model documents often contain mathematical expressions,
% definitions of the symbols used in those expressions, and ordinary
% textual terminology that should be available for downstream processing.
% The \textsf{measuredtex} package supports this by coupling ordinary
% \LaTeX\ markup with simple semantic annotations.
%
% The package is deliberately modest. It does not attempt to define a full
% model-interchange language. Instead, it records selected information from
% the source document in plain-text sidecar files during compilation. These
% sidecar files can then be inspected directly or processed by external
% tooling.
%
% The package provides three related facilities:
% \begin{itemize}
%   \item model environments that wrap standard \texttt{amsmath}
%     display environments and write model metadata;
%   \item commands and environments for recording definitions of
%     mathematical terms associated with a labelled model equation;
%   \item a command for marking textual concepts and writing a separate
%     concept list.
% \end{itemize}
%
% \section{Usage}
%
% Load the package in the usual way:
% \begin{verbatim}
%   \usepackage{measuredtex}
% \end{verbatim}
%
% \subsection{Model environments}
%
% \begin{function}{mdeq}
%   \begin{syntax}
%     \cs{begin}\{mdeq\}\oarg{metadata}
%       \meta{equation body}
%     \cs{end}\{mdeq\}
%   \end{syntax}
%
%   Typesets \meta{equation body} in a standard \texttt{equation}
%   environment and writes a model record to
%   \texttt{\cs{jobname}.models}.
%
%   The optional \meta{metadata} argument is a key--value list. The
%   recognised keys are:
%   \begin{description}
%     \item[\texttt{name}] a human-readable name for the model;
%     \item[\texttt{role}] the role of the equation in the document;
%     \item[\texttt{source}] a source or provenance string;
%     \item[\texttt{label}] the \LaTeX\ label to apply to the displayed
%       equation.
%   \end{description}
%
%   If the \texttt{label} key is supplied, the same value is written to the
%   sidecar file and applied as a \LaTeX\ label inside the equation
%   environment.
%
%   For example:
%   \begin{verbatim}
%     \begin{mdeq}[name=Ohm's law, role=model, label=eq:ohm]
%       V = IR
%     \end{mdeq}
%   \end{verbatim}
% \end{function}
%
% The body of the environment is typeset as an equation. A model record is
% written to \texttt{\cs{jobname}.models}. 
%
% \begin{function}{mdalign}
%   \begin{syntax}
%     \cs{begin}\{mdalign\}\oarg{metadata}
%       \meta{aligned equation body}
%     \cs{end}\{mdalign\}
%   \end{syntax}
%
%   The \texttt{mdalign} environment typesets
%   \meta{aligned equation body} in a standard \texttt{align}
%   environment and writes a model record to
%   \texttt{\cs{jobname}.models}.
%
%   The optional \meta{metadata} argument is a key--value list. The
%   recognised keys are \texttt{name}, \texttt{role}, \texttt{source},
%   and \texttt{label}.
%
%   If the \texttt{label} key is supplied, the same value is written to the
%   sidecar file and applied as a \LaTeX\ label inside the displayed
%   \texttt{align} environment.
%
%   For example:
%   \begin{verbatim}
%     \begin{mdalign}[name=Linear model, role=model, label=eq:linear]
%       y &= ax + b \\
%       u_y^2 &= x^2 u_a^2 + a^2 u_x^2 + u_b^2
%     \end{mdalign}
%   \end{verbatim}
% \end{function}
%
% \subsection{Mathematical term definitions}
%
% Mathematical terms may be associated with a model equation by referring
% to the equation label. Term records are written to
% \texttt{\cs{jobname}.models}. The package checks whether the supplied
% equation label has been defined and issues a warning if it has not.
%
% \begin{function}{\mdtext}
%   \begin{syntax}
%     \cs{mdtext}\marg{label}\marg{term}\oarg{separator}
%       \marg{definition}
%   \end{syntax}
%
%   Writes a mathematical-term definition record to
%   \texttt{\cs{jobname}.models} and also typesets an inline textual
%   definition.
%
%   The \meta{label} argument is the \LaTeX\ label of the model equation
%   to which the term belongs. The \meta{term} argument is typeset in
%   inline mathematics mode. The optional \meta{separator} argument is rendered
%   between the term and the definition text. 
%   The separator is not written to the sidecar file.
%   The \meta{definition} argument is rendered and written
%   to the sidecar file.
%
%   If \meta{label} has not been defined, the package issues an
%   \texttt{undefined-label} warning.
%
%   For example:
%   \begin{verbatim}
%     \mdtext{eq:ohm}{V}[:]{the voltage across the resistor}
%   \end{verbatim}
% \end{function}
%
% \begin{function}{\mdtab}
%   \begin{syntax}
%     \cs{mdtab}\marg{label}\marg{term}\marg{definition}
%   \end{syntax}
%
%   Writes a mathematical-term definition record to
%   \texttt{\cs{jobname}.models} and typesets a two-column table-row
%   fragment.
%
%   The \meta{label} argument is the \LaTeX\ label of the model equation
%   to which the term belongs. The \meta{term} argument is typeset in
%   inline mathematics mode. The rendered output consists of the term,
%   followed by an alignment character, followed by \meta{definition}.
%
%   This command is intended for use inside table-like environments such
%   as \texttt{tabular} or \texttt{tabularx}. It does not terminate the
%   row; the surrounding table markup should provide any required
%   \verb|\\|.
%
%   If \meta{label} has not been defined, the package issues an
%   \texttt{undefined-label} warning.
%
%   For example:
%   \begin{verbatim}
%     \begin{tabular}{ll}
%       \mdtab{eq:ohm}{V}{The voltage across the resistor} \\
%       \mdtab{eq:ohm}{I}{The current through the resistor} \\
%       \mdtab{eq:ohm}{R}{The resistance} \\
%     \end{tabular}
%   \end{verbatim}
% \end{function}
%
% \begin{function}{mddescription,\mditem}
%   \begin{syntax}
%     \cs{begin}\{mddescription\}
%       \cs{mditem}\marg{label}\marg{term}\marg{definition}
%       \ldots
%     \cs{end}\{mddescription\}
%   \end{syntax}
%
%   The \texttt{mddescription} environment provides a description-list
%   form for mathematical term definitions. It is a wrapper around the
%   standard \texttt{description} environment.
%
%   Within this environment, \cs{mditem}\marg{label}\marg{term}
%   \marg{definition} writes a mathematical-term definition record to
%   \texttt{\cs{jobname}.models} and typesets a description-list item.
%   The item label is \meta{term}, typeset in inline mathematics mode, and
%   the item body is \meta{definition}.
%
%   The \meta{label} argument is the \LaTeX\ label of the model equation
%   to which the term belongs. If \meta{label} has not been defined, the
%   package issues an \texttt{undefined-label} warning.
%
%   For example:
%   \begin{verbatim}
%     \begin{mddescription}
%       \mditem{eq:ohm}{V}{The voltage across the resistor}
%       \mditem{eq:ohm}{I}{The current through the resistor}
%       \mditem{eq:ohm}{R}{The resistance}
%     \end{mddescription}
%   \end{verbatim}
% \end{function}
%
% \begin{function}{\mddefn}
%   \begin{syntax}
%     \cs{mddefn}\marg{label}\marg{term}\marg{definition}
%   \end{syntax}
%
%   Writes a mathematical-term definition record to
%   \texttt{\cs{jobname}.models}. The \meta{label} argument is the
%   \LaTeX\ label of the model equation to which the term belongs.
%   The \meta{term} argument is the mathematical term being defined, and
%   \meta{definition} is the corresponding definition text.
%
%   This command does not typeset any visible material in the document.
%   It is intended for cases where the term definition should be recorded
%   semantically without being rendered at the point where the command
%   appears.
%
%   If \meta{label} has not been defined, the package issues an
%   \texttt{undefined-label} warning.
%
%   For example:
%   \begin{verbatim}
%     \mddefn{eq:ohm}{V}{The voltage across the resistor}
%   \end{verbatim}
% \end{function}
%
% \subsection{Textual concepts}
%
% \begin{function}{\mdconcept}
%   \begin{syntax}
%     \cs{mdconcept}\marg{id}\marg{text}
%   \end{syntax}
%
%   Marks an ordinary textual concept and records it for later output to the
%   concept sidecar file. The command typesets \meta{text} at the point
%   where it appears in the document, and records the pair
%   \meta{id}, \meta{text}.
%
%   Concept records are collected during the document and written 
%   to \texttt{\cs{jobname}.terms} at the  
%   end of document processing. Each distinct
%   concept ID is written once, in first-use order.
%
%   If an already defined \meta{id} is used with the same \meta{text}, the
%   existing concept record is reused silently. If the \meta{id} is
%   used again with different text, the package issues an
%   \texttt{inconsistent-concept} warning.
%
%   For example:
%   \begin{verbatim}
%     The \mdconcept{reference-oscillator}{reference oscillator}
%     provides the timing reference.
%   \end{verbatim}
%
%   The rendered document contains the text
%   \texttt{reference oscillator}. The concept sidecar file receives a
%   record of the form:
%   \begin{verbatim}
%     reference-oscillator: reference oscillator
%   \end{verbatim}
% \end{function}
%
% \subsection{Semantic term macros}
%
%   The \cs{mdterm} command generates mathematical terms according to 
%   well-defined semantic form.
%   The command is a formatting macro. It does not write
%   to a sidecar file.
%
% \begin{function}{\mdterm}
%   \begin{syntax}
%     \cs{mdterm}\oarg{namespace}\marg{symbol}\oarg{qualifier}
%       \texttt{(}\meta{association}\texttt{)}
%   \end{syntax}
%
%   Composes a decorated mathematical term from a
%   \meta{symbol} and optional namespace, qualifier, and association
%   components.
%
%   If \meta{namespace} is supplied, it is rendered using
%   \cs{nspace} (by default as a superscript). 
%   If \meta{qualifier} is supplied, it is rendered using \cs{qual}
%   (by default as a subscript). 
%   If the parenthesised \meta{association} is supplied, 
%   it is rendered after using \cs{assoc} (by default in parentheses after the symbol ).
%
%   For example:
%   \begin{verbatim}
%     \mdterm[SI]{V}[rms](input)
%   \end{verbatim}
% \end{function}
%
% The package uses small mathematical formatting macros to control the typesetting 
% of term notation. These macros may be redefined to change the format
% generated by \cs{mdterm}.
%
% \begin{function}{\nspace}
%   \begin{syntax}
%     \cs{nspace}\marg{text}
%   \end{syntax}
%
%   Typesets \meta{text} as a namespace component (default is upright roman
%   mathematics text).
%
%   For example:
%   \begin{verbatim}
%     \nspace{SI}
%   \end{verbatim}
% \end{function}
%
% \begin{function}{\assoc}
%   \begin{syntax}
%     \cs{assoc}\marg{text}
%   \end{syntax}
%
%   Typesets \meta{text} as an association component (default is upright roman
%   mathematics text).
%
%   For example:
%   \begin{verbatim}
%     \assoc{input}
%   \end{verbatim}
% \end{function}
%
% \begin{function}{\qual}
%   \begin{syntax}
%     \cs{qual}\marg{text}
%     \cs{qual}\marg{text}\texttt{(}\meta{sub-qualifier}\texttt{)}
%   \end{syntax}
%
%   Typesets \meta{text} as a qualifier (default is upright roman mathematics text
%   inside parentheses).
%
%   For example:
%   \begin{verbatim}
%     \qual{cal}
%     \qual{cal}(pre)
%   \end{verbatim}
% \end{function}
%
% \begin{function}{\tv}
%   \begin{syntax}
%     \cs{tv}\marg{text}
%   \end{syntax}
%
%   Typesets \meta{text} as a parametric term (default in italic mathematics text).
%
%   For example:
%   \begin{verbatim}
%     \tv{offset}
%   \end{verbatim}
% \end{function}
%
% \section{Sidecar file formats}
%
% \subsection{Model and term sidecar}
%
% The file \texttt{\cs{jobname}.models} is opened at the beginning of the
% document and closed at the end. Model records and mathematical-term
% definition records are written to this file.
%
% A model record begins with the header line:
% \begin{verbatim}
% ===== MODEL =====
% \end{verbatim}
%
% It may then contain any of the following metadata lines, depending on
% which keys were supplied:
% \begin{verbatim}
% name: <name>
% role: <role>
% source: <source>
% label: <label>
% equation: <equation-body>
% \end{verbatim}
%
% The \texttt{equation} line is always written for a model environment.
%
% Mathematical-term records are written as simple blocks containing:
% \begin{verbatim}
% ref: <equation-label>
% term: <term>
% definition: <definition>
% \end{verbatim}
%
% The current implementation does not add a separate term-record delimiter
% such as \texttt{----- TERM -----}.
%
% \subsection{Concept sidecar}
%
% The file \texttt{\cs{jobname}.terms} is written at the end of the
% document. It contains one block for each distinct concept ID recorded
% using \cs{mdconcept}. The output format is:
% \begin{verbatim}
% <id>: <text>
% \end{verbatim}
%
% Blank lines are written between concept records.
%
% \MaybeStop{%
%   \clearpage
%   \PrintChanges
%   \clearpage
%   \PrintIndex
% }
%
% \section{Implementation}
%
%    \begin{macrocode}
%<*package>
\NeedsTeXFormat{LaTeX2e}
\ProvidesPackage{measuredtex}[2026/06/05 v0.1 Semantic measurement markup]
\RequirePackage { amsmath }
\ExplSyntaxOn
%    \end{macrocode}
%
% \subsection{Messages}
%
% The package defines warning messages for undefined equation labels and
% for inconsistent reuse of textual concept identifiers.
%
%    \begin{macrocode}
\msg_new:nnn { measuredtex } { undefined-label }
  { md_defn_write:~label~'#1'~is~not~defined. }
%    \end{macrocode}
%
% \subsection{Metadata variables}
%%
%% These token-list variables hold the model metadata supplied through the
%% optional key--value argument of the model environments.
%%
%    \begin{macrocode}
\tl_new:N \l_md_name_tl
\tl_new:N \l_md_role_tl
\tl_new:N \l_md_source_tl
\tl_new:N \l_md_label_tl
\tl_new:N \l_md_model_tl
%    \end{macrocode}
%
% \subsection{Output stream}
%%
%% Two output streams are declared. The model stream is opened at the
%% beginning of the document and closed at the end. The terms stream is
%% opened only when the collected textual concepts are written.
%%
%    \begin{macrocode}
\iow_new:N \g_md_iow
\AtBeginDocument
  {
    \iow_open:Nn \g_md_iow { \jobname .models }
  }
\AtEndDocument
  {
    \iow_close:N \g_md_iow
  }
%    \end{macrocode}
% \subsection{Metadata keys}
%%
%% The \texttt{measuredtex/model} key family provides the supported model
%% metadata fields. The \texttt{measuredtex/mddefn} family currently
%% provides a \texttt{model} key, although the public term-definition
%% commands presently take positional arguments.
%%
%    \begin{macrocode}
\keys_define:nn { measuredtex / model }
  {
    name   .tl_set:N = \l_md_name_tl ,
    role   .tl_set:N = \l_md_role_tl ,
    source .tl_set:N = \l_md_source_tl ,
    label  .tl_set:N = \l_md_label_tl ,
  }
\keys_define:nn { measuredtex / mddefn }
  {
    model .tl_set:N = \l_md_model_tl ,
  }
%    \end{macrocode}
% \subsection{Model-writing logic}
%%
%% \begin{macro}[int]{\md_model_write:nn}
%% Clears all metadata variables, applies the supplied key--value pairs,
%% and writes a model block to the sidecar file. Empty optional metadata
%% fields are omitted. The equation body is written in string form.
%%
%    \begin{macrocode}
\cs_new:Npn \md_model_write:nn #1 #2
  {
    \tl_clear:N \l_md_name_tl
    \tl_clear:N \l_md_role_tl
    \tl_clear:N \l_md_source_tl
    \tl_clear:N \l_md_label_tl
    \keys_set:nn { measuredtex / model } { #1 }
    \iow_now:Nn \g_md_iow { ===== ~ MODEL ~ ===== }
    \tl_if_empty:NF \l_md_name_tl
      {
        \iow_now:Nx \g_md_iow { name: ~ \tl_use:N \l_md_name_tl }
      }
    \tl_if_empty:NF \l_md_role_tl
      {
        \iow_now:Nx \g_md_iow { role: ~ \tl_use:N \l_md_role_tl }
      }
    \tl_if_empty:NF \l_md_source_tl
      {
        \iow_now:Nx \g_md_iow { source: ~ \tl_use:N \l_md_source_tl }
      }
    \tl_if_empty:NF \l_md_label_tl
      {
        \iow_now:Nx \g_md_iow { label: ~ \tl_use:N \l_md_label_tl }
      }
    \iow_now:Nx \g_md_iow { equation: ~ \tl_to_str:n {#2} }
    \iow_now:Nn \g_md_iow { }
  }
%    \end{macrocode}
% \end{macro}
%
% \subsection{Document environments}
% \begin{macro}{mdeq}
%%
%% The \texttt{mdeq} environment records the model metadata and body, then
%% typesets the body inside an \texttt{equation} environment. If a label was
%% supplied through the metadata keys, it is applied inside the equation.
%%
%    \begin{macrocode}
\NewDocumentEnvironment { mdeq } { O{} +b }
  {
    \md_model_write:nn {#1} {#2}
    \begin{equation}
      #2
      \tl_if_empty:NF \l_md_label_tl
        {
          \label { \tl_use:N \l_md_label_tl }
        }
    \end{equation}
  }
  { }
%    \end{macrocode}
% \end{macro}
%%
% \begin{macro}{mdalign}
%%
%% The \texttt{mdalign} environment is the corresponding wrapper for
%% \texttt{align}.
%    \begin{macrocode}
\NewDocumentEnvironment { mdalign } { O{} +b }
  {
    \md_model_write:nn {#1} {#2}
    \begin{align}
      #2
      \tl_if_empty:NF \l_md_label_tl
        {
          \label { \tl_use:N \l_md_label_tl }
        }
    \end{align}
  }
  { }
%    \end{macrocode}
% \end{macro}
% \subsection{Term definition commands}
%%
% \begin{macro}[int]{\md_defn_write:nnn}
%% The internal term-writing function checks that the referenced equation
%% label exists and writes the reference, term, and definition fields to the
%% model sidecar file.
%    \begin{macrocode}
\bool_new:N \g_md_suppress_nested_sidecar_bool

\cs_new_protected:Npn \md_defn_write:nnn #1 #2 #3
  {
    \cs_if_exist:cF { r@#1 }
      {
        \msg_warning:nnn { measuredtex } { undefined-label } {#1}
      }
    \iow_now:Nn  \g_md_iow { type: ~ definition }
    \iow_now:Nx  \g_md_iow { ref: ~ \tl_to_str:n {#1} }
    \iow_now:Nx  \g_md_iow { term: ~ \tl_to_str:n {#2} }
    \iow_now:Nx  \g_md_iow { description: ~ \tl_to_str:n {#3} }
    \iow_now:Nn  \g_md_iow { }
  }
\NewDocumentCommand \mddefn { m m m }
  {
    \md_defn_write:nnn {#1} {#2} {#3}
  }
\NewDocumentCommand \mdtext { m m O{} m }
  {
    \md_defn_write:nnn {#1} {#2} {#4}
    \group_begin:
      \bool_set_true:N \g_md_suppress_nested_sidecar_bool
      $#2$#3~#4
    \group_end:
  }    
\NewDocumentCommand \mdtab { m m m }
  {
    \md_defn_write:nnn {#1} {#2} {#3}
    \bool_gset_true:N \g_md_suppress_nested_sidecar_bool
    $#2$ & #3
    \bool_gset_false:N \g_md_suppress_nested_sidecar_bool
  }
\RequirePackage { enumitem  }
\NewDocumentEnvironment { mddescription } { O{} }
  {
    \begin{description}[labelindent=\leftmargini,#1]
  }
  {
    \end{description}
  }
\NewDocumentCommand \mditem { m m m }
  {
    \md_defn_write:nnn {#1} {#2} {#3}
    \group_begin:
      \bool_set_true:N \g_md_suppress_nested_sidecar_bool
      \item[${#2}$] #3
    \group_end:
  }   
\prop_new:N \g_md_concepts_prop
\seq_new:N  \g_md_concepts_seq
\msg_new:nnn { measuredtex } { inconsistent-concept }
  {
    Concept~ID~'#1'~already~defined~as~'#2',~not~'#3'.
  }    
\tl_new:N \l_md_concept_id_tl

\cs_new_protected:Npn \md_make_concept_id:nN #1 #2
  {
    \tl_set:Nn #2 {#1}
    \tl_trim_spaces:N #2
    \regex_replace_all:nnN { \s+ } { _ } #2
  }

\NewDocumentCommand \mdconcept { o m }
  {
    \IfValueTF {#1}
      { \tl_set:Nn \l_md_concept_id_tl {#1} }
      { \md_make_concept_id:nN {#2} \l_md_concept_id_tl }

    #2

    \prop_if_in:NVTF \g_md_concepts_prop \l_md_concept_id_tl
      {
        \prop_get:NVN \g_md_concepts_prop \l_md_concept_id_tl \l_tmpa_tl
        \tl_if_eq:NnF \l_tmpa_tl {#2}
          {
            \msg_warning:nnnnn
              { measuredtex }
              { inconsistent-concept }
              { \l_md_concept_id_tl }
              { \l_tmpa_tl }
              {#2}
          }
      }
      {
        \prop_gput:NVn \g_md_concepts_prop \l_md_concept_id_tl {#2}
        \seq_gput_right:NV \g_md_concepts_seq \l_md_concept_id_tl
      }
      
    % Write to sidecar unless suppressed by an outer wrapper    
    \bool_if:NF \g_md_suppress_nested_sidecar_bool
      {
        \iow_now:Nn \g_md_iow { type: ~ concept }
        \iow_now:Nx \g_md_iow
          {
            markup: ~
            \c_backslash_str mdconcept ~
            [
            \tl_to_str:V \l_md_concept_id_tl
            ]
            {
            \tl_to_str:n {#2}
            }
          }
        \iow_now:Nn \g_md_iow { }
      }
  }
%    \end{macrocode}
% \end{macro}
%
% \subsection{Semantic term macros}
%%
%% These commands use conventional \LaTeXe\ syntax, so the implementation
%% temporarily leaves \texttt{expl3} syntax.
%%
%    \begin{macrocode}
\ExplSyntaxOff
%    \end{macrocode}
%%
% \begin{macro}{\nspace}
%% Renders a namespace component in upright roman maths text.
%    \begin{macrocode}
\NewDocumentCommand{\nspace}{m}{%
  \mathrm{#1}%
}
%    \end{macrocode}
% \end{macro}
%%
% \begin{macro}{\assoc}
%% Renders an association component in upright roman maths text.
%    \begin{macrocode}
\NewDocumentCommand{\assoc}{m}{%
  \mathrm{#1}%
}
%    \end{macrocode}
% \end{macro}
%%
% \begin{macro}{\tv}
%% Renders a term value in italic maths text.
%    \begin{macrocode}
\NewDocumentCommand{\tv}{m}{%
  \mathit{#1}%
}
%    \end{macrocode}
% \end{macro}
%%
% \begin{macro}{\qual}
%% Renders a qualifier, with an optional parenthesised sub-qualifier.
%    \begin{macrocode}
\NewDocumentCommand{\qual}{m d()}{%
  \mathrm{#1}\IfValueT{#2}{(\mathrm{#2})}%
}
%    \end{macrocode}
% \end{macro}
%%
% \begin{macro}{\mdterm}
%% Composes a semantic mathematical term from an optional namespace, a
%% required symbol, an optional qualifier, and an optional association.
%    \begin{macrocode}
\NewDocumentCommand{\mdterm}{o m o d()}{%
  \IfValueT{#1}%
    {{}^{\nspace{#1}}}%
  {#2}%
  \IfValueT{#3}%
    {_{\qual{#3}}}%
  \IfValueT{#4}%
    {(\assoc{#4})}%
}
%    \end{macrocode}
% \end{macro}
%    \begin{macrocode}
\endinput
%</package>
%    \end{macrocode}
%
% \Finale